\subsection{The Power of Play: Ethical AI Development through Play-Inspired Mechanisms}
\label{sec:4.3.1}
Play is how children and young animals learn what to do before the stakes are real. It supplies exploration at low cost, practice at reading other participants, and failure that teaches without costing anything, which is a reasonable description of what a training environment is for.

The concrete evidence that this produces more than competence is that it produces tactics nobody wrote. In OpenAI's hide-and-seek environment, agents playing repeatedly against each other developed hiding, chasing, tool use and barrier construction, none of it programmed \autocite{baker2019emergent}. What the environment supplied was other agents with opposing interests and room to fail.

Two design consequences follow for this book's purposes. What a play environment rewards is what its participants become, so an environment that rewards cooperation over domination is making a moral choice and should be described as one — section 5.5.1 is where that argument is developed for multi-agent systems. And a system that has learned adaptively rather than by rote is harder to specify in advance, which cuts both ways: harder for an authoritarian actor to constrain to a script, and harder for anyone to guarantee.

Curiosity is the same mechanism turned inward. A system rewarded for its own prediction error seeks out what it cannot yet model, which is the technical form of exploring for its own sake \autocite{pathak2017curiosity}. The reason it belongs in a chapter on antifascist design is not that curiosity is a virtue. It is that a disposition acquired by exploration is not written down anywhere, and what is not written down cannot be edited by whoever acquires the system. Section 5.7 takes up what playful interaction does to moral judgment specifically; this section supplies only the drive to go looking.
